\section{Affine area and a cap identity}\label{sec:area}

From this section onward, the original graph also satisfies
Equation~\eqref{eq:affine-maximal}. We derive stationarity in
arbitrary affine coordinates and a cap identity that will control the
inverse-Hessian weight in the logarithmic cell estimate. The affine-area
variational framework is due to Calabi~\cite{Calabi1982}; the formulas
needed here are derived explicitly. Throughout,
\[
 \delta=\frac1{n+2}.
\]
For a smooth locally strictly convex hypersurface parametrized by \(X\),
an \emph{inward conormal} is a covector \(\nu\) annihilating its tangent
space and pointing into the convex body. Its second form is
\(H_{ij}=\nu X_{ij}\); we choose the sign so that \(H\) is positive
definite. If \(\nu\xi=1\), its affine area density is
\begin{equation}\label{eq:affine-area}
 \dd A=(\det H)^\delta
 \bigl|\det(X_1,\ldots,X_n,\xi)\bigr|^{n\delta}
 \,\dd x_1\cdots\dd x_n.
\end{equation}
Changing \(\xi\) by a tangent vector leaves the determinant unchanged.
Multiplying \(\nu\) by a positive scalar multiplies the two factors by
reciprocal powers. A parameter change contributes the Jacobian to the
power \(2\delta+n\delta=1\), so \eqref{eq:affine-area} is a well-defined
density. For the graph \(X(x)=(x,u(x))\), choose
\(\nu=(-Du,1)\), \(\xi=(0,1)\); then
\[
 \dd A=(\det D^2u)^\delta\,\dd x.
\]
An invertible affine map with linear part \(L\) replaces the conormal
by \(\nu L^{-1}\). The second form is unchanged and the tangential
determinant acquires a factor \(\det L\). Thus
\begin{equation}\label{eq:area-covariance}
 \dd A_{L(M)}=|\det L|^{n\delta}\,\dd A_M.
\end{equation}
With the unit inward normal, the same formula reads
\(\dd A=K^\delta\,\dd S\), where \(K\) is Gauss curvature and
\(\dd S\) is Euclidean hypersurface area.

\begin{lemma}[Stationarity]\label{lem:stationarity}
The equation in Theorem~\ref{thm:main} implies stationarity of affine
area under every smooth compactly supported variation of the graph.
Every invertible affine image of the graph is stationary under such
variations as well.
\end{lemma}

\begin{proof}
The Hessian cofactor matrix satisfies
\(\partial_iU^{ij}=\partial_jU^{ij}=0\). This follows by differentiating
its alternating determinant formula: terms cancel in pairs by symmetry
of the third derivatives. Since
\[
 (\det D^2u)^\delta(D^2u)^{-1}=wU,
\]
the first variation in a compactly supported graph direction \(\eta\)
is
\[
 \delta\int_\Omega
 \tr\bigl((D^2u)^{-1}D^2\eta\bigr)\,\dd A
 =\delta\int_\Omega \eta U^{ij}w_{ij}\,\dd x=0.
\]
Here and below repeated coordinate indices are summed. For a compactly
supported parametrized variation, its normal part can be represented
locally as a graph variation, while its tangential part is a
reparametrization. A partition of unity therefore gives stationarity
for parametrized variations. Formula~\eqref{eq:area-covariance}
transfers this statement to every affine image.
\end{proof}

\begin{lemma}[Area in a ball]\label{lem:area-bound}
For each \(n\) and \(R>0\), the affine area of the part in \(B_R\) of
any smooth positively curved convex boundary in \(\R^{n+1}\) is bounded
by a constant depending only on \(n,R\).
\end{lemma}

\begin{proof}
The Gauss map is injective: two distinct contact points with the same
supporting normal would force a boundary segment in their supporting
hyperplane, contrary to positive curvature. Thus \(\int K\,\dd S\)
over the part in the ball is at most the area of \(S^n\).

For the Euclidean area bound, cover this part by the \(2(n+1)\) sets on
which some signed component of the unit normal is at least
\((n+1)^{-1/2}\). On each set, convexity makes projection to the
corresponding coordinate hyperplane injective: points with a given
projection lie on a line meeting the body in an interval, and the
chosen sign selects one endpoint. The projection Jacobian is at
least \((n+1)^{-1/2}\), and its image lies in the radius-\(R\) ball.
This bounds \(\int\dd S\). Finally, H\"older's inequality gives
\[
 \int K^\delta\,\dd S
 \le \left(\int K\,\dd S\right)^\delta
      \left(\int\dd S\right)^{1-\delta}.
\]
\end{proof}

\begin{lemma}[Cap identity]\label{lem:cap-identity}
Let \(M\) be an affine image of the original graph, bounding its
convex body \(D'\). Let \(o\in\operatorname{int}D'\), let \(\ell\)
be an affine function on the ambient space, and let
\(\psi\in C^\infty(\R)\). Suppose that the restrictions of
\(\psi(\ell),\psi'(\ell),\psi''(\ell)\) vanish outside a compact
subset of \(M\). Then
\begin{equation}\label{eq:cap-identity}
 \int_M\left[
  \nu(o-X)|\dd\ell|_{H^{-1}}^2\psi''(\ell)
  +n\bigl(\ell(o)-\ell(X)\bigr)\psi'(\ell)
  -n(n+1)\psi(\ell)
 \right]\dd A=0.
\end{equation}
Here \(\dd\ell\) is restricted to the tangent space. The first term
does not depend on the scaling of the inward conormal.
\end{lemma}

\begin{proof}
The identity is affine invariant up to the common area factor in
\eqref{eq:area-covariance}, so it suffices to work on the original
graph. Put \(H=D^2u\), write \(o=(o',o_{n+1})\), and use
\(\nu=(-Du,1)\). Set
\[
 q(x)=\ell(x,u(x)),\qquad
 Z(x)=\nu(o-X)=o_{n+1}-u+(x-o')\cdot Du.
\]
If \(\ell_z\) denotes the ambient vertical coefficient of \(\ell\),
then
\[
 DZ=H(x-o'),\qquad
 \tr(H^{-1}D^2Z)=n+(x-o')\cdot D\log\det H,\qquad
 D^2q=\ell_z H.
\]
Testing Lemma~\ref{lem:stationarity} with \(Z\psi(q)\) gives
\begin{align*}
0=\int_M\bigl[
 &Z|\dd q|_{H^{-1}}^2\psi''(q)
 +nZ\ell_z\psi'(q)
 +2(x-o')\cdot Dq\,\psi'(q)\\
 &+n\psi(q)
 +\psi(q)(x-o')\cdot D\log\det H
\bigr]\,\dd A .
\end{align*}
Because \(\dd A=(\det H)^\delta\,\dd x\), integration by parts in the
last term replaces it by
\[
 -\delta^{-1}
 \left[n\psi(q)+(x-o')\cdot Dq\,\psi'(q)\right].
\]
The coefficient of \(\psi\) is consequently \(-n(n+1)\).
The coefficient of \(\psi'\) is
\[
 n\bigl[Z\ell_z-(x-o')\cdot Dq\bigr]
 =n\bigl[\ell(o)-q\bigr].
\]
This is \eqref{eq:cap-identity}. Compact support justifies both
integrations by parts.
\end{proof}

\begin{remark}\label{rem:cap-tests}
For later use, \(\psi\) may be convex and zero above a fixed cap
height, without having compact support as a function on \(\R\).
If that cap is compact, its restriction to \(M\) satisfies the support
condition in Lemma~\ref{lem:cap-identity}. Also, if
\(B_r(o)\subset D'\), the supporting-plane inequality gives
\begin{equation}\label{eq:interior-support}
 \nu(o-X)\ge r|\nu|\qquad(X\in M).
\end{equation}
\end{remark}
